\begin{statement}{B.14d}{Proved}{Uniform integral symmetry lifting at both exceptional primes}
For $p=2,3$ and every $n\geq1$, let
$T_{n,p}=\mathfrak s_{2,n}\otimes\mathbb Z_p$ with its
geometric $S_3$-action. Then
\[
 H^1(S_3,T_{n,p})=0.
\]
Let $P_n$ be the degree-$n$ subgroup of special outer
IA automorphisms of the free pro-$p$ tripod group, as in
B.12. The map
\[
 P_n^{S_3}\longrightarrow(P_n/P_{n+1})^{S_3}
\]
is surjective for every $n$. Thus every integral
$S_3$-invariant special outer leading term lifts to an
actual continuous special outer automorphism commuting
with the geometric $S_3$. In particular $q_7$, $2q_7$,
and $3q_7$ have infinite pro-$2$ and pro-$3$ lifts without
any additional multiplier. Together with B.12 this gives
the assertion for every prime. For each prime the lift
of any of these three terms satisfies the two cycles
and fails the pentagon.
\end{statement}
\begin{proof}
\emph{The actual integral module.}
Put $R=\mathbb Z_p$, $V=R^3/R(1,1,1)$, and $L_m=L_m(V)$.
In the basis $x,y$, with $z=-x-y$, the generators of $S_3$
act by
\[
 \tau(x)=y,\quad\tau(y)=x,
 \qquad \rho(x)=y,\quad\rho(y)=-x-y.
\]
For $n\geq2$, the equivariant sequence of B.10 is
\begin{equation}\label{eq:B14d-special}
 0\longrightarrow T_{n,p}\longrightarrow
 D_n:=V\otimes_R L_n\xrightarrow{b}L_{n+1}
 \longrightarrow0.
\end{equation}
Its integral surjectivity was proved by Jacobi in B.10.
The identification uses the three peripheral conjugators
modulo their common inner conjugator; in degree $n\geq2$
their generator-centralizer ambiguities are zero. Thus
there is no further quotient to take after forming this
kernel. It is the same module, with the same tensor
action, used in B.13 and B.14c. The degree-one module
is zero, and the bracket map in degree two is an
isomorphism, so $T_{1,p}=T_{2,p}=0$.

\emph{The prime two.}
Let $C_2=\langle\tau\rangle$ and write $N_2=1+\tau$.
The integral basis theorem of
\citet[Theorem 2, pp.~15--16]{BryantKovacsStohr2002InvariantBases}
applies to the free Lie ring on the swapped pair $x,y$.
It gives a homogeneous integral basis consisting of
swapped pairs and sign eigenvectors, with no trivial
one-dimensional summands. On a swapped pair the invariant
vectors are exactly the norms; a sign line has no
invariant vector over the torsion-free ring $R$.
Consequently, in every degree,
\begin{equation}\label{eq:B14d-norm2}
 L_m^{C_2}=N_2L_m.
\end{equation}

The restriction of $V$ to $C_2$ is the regular lattice
$R[C_2]$. Its tensor product with any finite free
$R$-lattice, with diagonal action, is again a free
$R[C_2]$-module: on a group-basis tensor the change
$g\otimes v\mapsto g\otimes g^{-1}v$ removes the action
on the second factor. Hence $H^1(C_2,D_n)=0$.
Moreover $b:D_n^{C_2}\to L_{n+1}^{C_2}$ is onto:
given $c=N_2w$, lift $w$ to $u\in D_n$ by the bracket
surjectivity, and use $N_2u$ as an invariant lift.
The cohomology sequence of \eqref{eq:B14d-special}
therefore gives $H^1(C_2,T_{n,2})=0$.

To pass to $S_3$, its normal subgroup $C_3$ has order
invertible in $\mathbb Z_2$. The operator
$e_3=(1+\rho+\rho^2)/3$ is an $S_3$-equivariant
projection onto $T_{n,2}^{C_3}$. Thus the latter is a
$C_2$-equivariant direct summand, and its $C_2$ first
cohomology vanishes. Explicitly, any $S_3$ cocycle can
first be changed by a coboundary to vanish on $C_3$, by
averaging there. Its remaining values form a $C_2$
cocycle in $T_{n,2}^{C_3}$; this is a coboundary as well.
This proves $H^1(S_3,T_{n,2})=0$ in every degree.

\emph{The prime three.}
Let $C_3=\langle\rho\rangle$, $N_3=1+\rho+\rho^2$,
and $A=R[C_3]/RN_3$. The displayed action identifies
$V$ with $A$; multiplication by $\rho-1$ also identifies
$A$ with the augmentation ideal. In particular $A^{C_3}=0$.

We use the integral decomposition theorem of
\citet[Theorem 7.1, p.~453]{KovacsStohr2000LiePowers}.
For the exact $S_3$ generated by the displayed substitutions,
every integral $L_m$ with $m\ne2$ is a direct sum of
copies of $L_1,L_3,L_4,L_5$. Their paper uses row-vector
right actions; transposing its matrices gives exactly
our column-vector conventions.
On restriction to $C_3$ over $R$, both $L_1$ and $L_3$
are $A$: integrally the map
\[
 V\otimes\bigwedge^2 V\longrightarrow L_3,
 \qquad v\otimes(x\wedge y)\longmapsto[v,[x,y]]
\]
is an isomorphism, and $\rho$ has determinant one.
The restrictions of $L_4$ and $L_5$ are respectively
free $R[C_3]$-modules of ranks one and two. To verify
this integral assertion, their reductions modulo $3$
are projective by
\citet[Corollary 4.4, p.~441]{KovacsStohr2000LiePowers}.
Restriction to $C_3$ is then free over the local ring
$\mathbb F_3[C_3]$. Lift a modular free basis to $L_m$.
The resulting $R[C_3]$-map from a free module is square
as an $R$-linear map and has determinant a unit modulo
$3$, so it is an isomorphism over $R$.

It follows that, for every $m\ne2$, $L_m$ restricted
to $C_3$ is a direct sum of augmentation and regular
lattices. Both have invariant vectors equal to norms:
the augmentation lattice has no invariant vectors,
while the regular lattice has invariant line $RN_3$.
Thus
\begin{equation}\label{eq:B14d-norm3}
 L_m^{C_3}=N_3L_m\qquad(m\ne2).
\end{equation}
The elementary periodic resolution of a cyclic group
gives $H^2(C_3,L_m)=L_m^{C_3}/N_3L_m=0$ in these degrees.

Tensor the exact sequence
$0\to R\xrightarrow{1\mapsto N_3}R[C_3]\to V\to0$
with $L_n$, keeping the diagonal action. Its middle term
is free over $R[C_3]$ by the same untwisting of tensor
factors as above, and hence has zero positive cohomology.
For $n\geq3$ its long exact sequence gives
\[
 H^1(C_3,D_n)\simeq H^2(C_3,L_n)=0.
\]
Also $b:D_n^{C_3}\to L_{n+1}^{C_3}$ is onto:
write an invariant vector as $N_3w$ using
\eqref{eq:B14d-norm3}, lift $w$ through $b$, and take
its norm. Equation \eqref{eq:B14d-special} now gives
$H^1(C_3,T_{n,3})=0$ for every $n\geq3$; the two
smaller degrees are already zero.
An $S_3$ cocycle can therefore first be corrected to
vanish on $C_3$. It then factors through a $C_2$ cocycle
in $T_{n,3}^{C_3}$, which vanishes by averaging because
$2$ is a unit in $R$. This proves the required uniform
vanishing at $p=3$.

\emph{The uniform passage to actual automorphisms.}
B.11 identifies $P_m/P_{m+1}$ with $T_{m,p}$ and
supplies a lift in $P_m$ of every integral vector in
that module. Let $v\in T_{n,p}^{S_3}$, and choose an
unrestricted lift $a\in P_n$. Its discrepancies
$a^{-1}g(a)$ lie in $P_{n+1}$.
Suppose inductively that they lie in $P_m$.
Their classes $c(g)$ in $T_{m,p}$ satisfy
\[
 c(gh)=c(g)+g\,c(h),
\]
as follows by expanding
$a^{-1}gh(a)=(a^{-1}g(a))\,g(a^{-1}h(a))$.
Uniform first-cohomology vanishing supplies
$b\in T_{m,p}$ with $c(g)=b-gb$.
Choose $u\in P_m$ lifting $b$. Replacing $a$ by $au$
changes the discrepancy class to
$-b+c(g)+gb=0$. It therefore improves the symmetry
by one degree, without changing its prescribed leading
term. This procedure applies in every subsequent degree.

The successive corrections converge in the actual
pro-$p$ automorphism topology. Indeed every finite
characteristic quotient of a finitely generated free
pro-$p$ group is a nilpotent finite $p$-group, so a
sufficiently deep lower-central correction acts trivially
on it. The special outer IA automorphism group is
compact: it is a closed subgroup of the profinite outer
automorphism group, the peripheral conjugacy conditions
being closed. Its degree filtration is separated. For
the latter assertion, an automorphism that is inner
modulo every lower-central quotient is inner globally:
the nonempty closed sets of conjugating elements in
the compact free pro-$p$ group have the finite-intersection
property. The limit of the corrected outer automorphisms
therefore commutes with every element of $S_3$ and
still has leading term $v$.

This is a uniform argument on the actual group; it uses
neither an exponential with uncontrolled denominators
nor an extrapolation from the class-fourteen calculation.
The normalized two-cycle/three-cycle identification of
B.12 uses only free-group substitutions and peripheral
centralizers, and applies at these primes as well.
Finally the first pentagon defect of $m q_7$ has the
coefficient $16m$ recorded in B.12. For $m=1,2,3$ this
is nonzero in every $\mathbb Z_p$. Higher-degree
corrections do not alter that first defect. The lifts
therefore satisfy the two cycles and fail the pentagon,
as asserted. This conclusion is prime by prime and
does not assemble the lifts into a full profinite one.
\end{proof}
